%% -!TEX root = AUTthesis.tex
% در این فایل، عنوان پایان‌نامه، مشخصات خود، متن تقدیمی‌، ستایش، سپاس‌گزاری و چکیده پایان‌نامه را به فارسی، وارد کنید.
% توجه داشته باشید که جدول حاوی مشخصات پروژه/پایان‌نامه/رساله و همچنین، مشخصات داخل آن، به طور خودکار، درج می‌شود.
%%%%%%%%%%%%%%%%%%%%%%%%%%%%%%%%%%%%
% دانشکده، آموزشکده و یا پژوهشکده  خود را وارد کنید
\faculty{دانشکده مهندسی برق}
% گرایش و گروه آموزشی خود را وارد کنید
\department{گرایش کنترل}
% عنوان پایان‌نامه را وارد کنید
\fatitle{طراحی مسیر و پیاده‌سازی کنترل پلتون  خودروهای  
\\[.75 cm]
 هوشمند در بستر اینترنت اشیاء}
% نام استاد(ان) راهنما را وارد کنید
\firstsupervisor{دکتر حیدرعلی طالبی}
\secondsupervisor{دکتر ایمان شریفی}
% نام استاد(دان) مشاور را وارد کنید. چنانچه استاد مشاور ندارید، دستور پایین را غیرفعال کنید.
\firstadvisor{دکتر حیدرعلی طالبی}
%\secondadvisor{دکتر ایمان شریفی}
% نام نویسنده را وارد کنید
\name{روزبه }
% نام خانوادگی نویسنده را وارد کنید
\surname{بازرگانی}
%%%%%%%%%%%%%%%%%%%%%%%%%%%%%%%%%%
\thesisdate{فروردین 1400}

% چکیده پایان‌نامه را وارد کنید
\fa-abstract{
در اين قسمت چكيده پایان نامه نوشته مي‌شو‌د‌.‌ چكيده بايد جامع و بيان‌كننده‌ خلاصه‌اي از اقدامات انجام‌شده باشد. در چكيده باید از ارجاع به مرجع و ذكر روابط رياضي، بيان تاريخچه و تعريف مسئله خودداري ‌شود. 
}


% کلمات کلیدی پایان‌نامه را وارد کنید
\keywords{کلیدواژه اول، ...، کلیدواژه پنجم (نوشتن سه تا پنج واژه کلیدی ضروری است)}



\AUTtitle
%%%%%%%%%%%%%%%%%%%%%%%%%%%%%%%%%%
\vspace*{7cm}
\thispagestyle{empty}
\begin{center}
\includegraphics[height=5cm,width=12cm]{besm}
\end{center}